\chapter{Anti-GBM Antibodies}

\section*{What Is This Test?}
Anti-glomerular basement membrane antibodies target components of the basement membrane in the kidney and lung. They are an important serologic marker of anti-GBM disease, which can cause rapidly progressive glomerulonephritis with or without pulmonary hemorrhage.

\section*{Alternative Names}
Anti-GBM; Glomerular Basement Membrane Antibody; Goodpasture Antibody.

\section*{Why This Test Is Ordered}
Testing is urgent when rapidly worsening kidney function, blood or protein in the urine, pulmonary hemorrhage, or a pulmonary-renal syndrome raises concern for anti-GBM disease. It is usually ordered with ANCA, ANA, complement, urinalysis, and kidney-function tests.

\section*{How This Test Is Performed}
A blood sample is tested by an immunoassay for antibodies against the alpha-3 chain of type IV collagen. Kidney or lung tissue may sometimes be examined, but the blood test is the focus of this chapter.

\section*{Preparation, Risks, and Limitations}
No fasting is required. Venipuncture is the usual risk. A negative result does not exclude every case, and a positive result must be interpreted with urinalysis, kidney function, symptoms, and specialist evaluation.

\section*{Understanding Results}
A strong positive result in a compatible pulmonary-renal syndrome is highly concerning for anti-GBM disease and requires urgent specialist management. Low-level or unexpected positivity requires confirmation and assessment for other autoimmune or inflammatory conditions.

\section*{References}
Kidney Disease: Improving Global Outcomes guidance on rapidly progressive glomerulonephritis.

\clearpage
\section*{Understanding Results and Follow-up}
The laboratory report should identify the specimen, method, units, reference interval or decision limit, and whether the result is positive, negative, abnormal, or inconclusive. Reference intervals vary with age, sex, pregnancy, specimen type, assay, and laboratory; a result outside a range is not automatically a diagnosis, and a result within a range does not always exclude disease. Discuss unexpected results with the ordering clinician, who can decide whether repeat or confirmatory testing, treatment, or referral is appropriate. Do not change medicines or delay urgent care based on an isolated result.


\section*{Regulatory Status and Authorization Date}
Regulatory status applies to the specific assay, instrument, manufacturer, and intended use rather than to the test name alone. No single approval date can be assigned to this test category. For a commercial assay, verify the current product in the relevant regulator's database: the U.S. Food and Drug Administration (FDA) clearance, approval, or authorization; India's Central Drugs Standard Control Organisation (CDSCO) licence or permission; the European Union In Vitro Diagnostic Regulation (EU IVDR) CE certificate issued through the applicable conformity-assessment route; or the United Kingdom Medicines and Healthcare products Regulatory Agency (MHRA) registration and UKCA/CE status. Laboratory-developed tests may not have an individual FDA, CDSCO, EU IVDR, or MHRA approval date.
\clearpage
